\section{Rational approximation and coefficient height}
\label{c-approximation}

This appendix follows source Appendix~C, pp.~28--31, including all steps
of its (C.1)--(C.8).  The matrix perturbation calculation survives the
change of moment order.  We first derive its relative norm estimate,
then show exactly where a negative completion budget would be required.
Throughout, $K=40n$, $N=3n$, $h=37n$, $\alpha=3/40$, and $\lambda=37/40$.

\subsection{The relative matrix norm}
\label{c-relative-norm}

Write $\xi=\zeta(9)$, $G=G_{9,K}(\xi)$, and
$M_{9,K}=[X]G_{9,K}(X)$.  The positive Gram representation makes $G$
positive definite.  Consequently the real symmetric matrix
$B_K=G^{-1/2}M_{9,K}G^{-1/2}$ is defined.  Its Euclidean operator norm
measures a change of $X$ relative to the size of the unperturbed matrix.
The analogue of source (C.1) is
\begin{equation}
 \|B_K\|\leq e^{\rho K},\qquad
 \rho=\frac{6933}{500}=13.866,
 \quad K\in40\mathbb Z_{>0}\text{ sufficiently large}.
 \label{c-norm-bound}
\end{equation}
The value of $\rho$ is the source's fraction $6933/500$; its decimal is
$13.866$.

\noindent\textbf{[Derived]} Let $q$ be a nonzero real polynomial of
degree less than $h$, write $\mathbf q$ for its coefficient vector, and
put $p(t)=q(K^2t)$.  For fixed $0<a<b$, retain just the $\ell=1$ term in
the positive weight
\[
 w_9(y)=\frac{(2\pi)^8y^9}{20160}
          \sum_{\ell\geq1}\ell^8e^{-2\pi\ell y}.
\]
With $y=K\sqrt t$, this term contributes
$K^{10}(2\pi)^8t^4e^{-2\pi K\sqrt t}\,dt/40320$.
On $[a,b]$ the scaled numerator is at least $a^{6N}$ and the scaled
denominator at most $(1+b)^K$.  The analogue of source (C.2) is therefore
\begin{equation}
 \mathbf q^TG\mathbf q\geq
 A_Kc_a e^{-F(a,b)K}\int_a^b p(t)^2\,dt,
 \quad
 \begin{aligned}
 A_K&=K^{12N-2K+10},&
 c_a&=\frac{(2\pi)^8a^4}{40320},\\
 F(a,b)&=2\pi\sqrt b+\log(1+b)-6\alpha\log a.
 \end{aligned}
 \label{c-gram-lower}
\end{equation}
The extra four powers of $K$ relative to the source occur in both
quadratic forms below, so they cancel from their ratio.
Restricting the positive integral is useful even though it discards
most of the measure: it controls the entire coefficient vector through
the squared norm of $p$ on one fixed interval.  No cancellation enters
this lower bound.  The later Legendre estimate converts that interval
norm into control of the evaluations at the negative pole locations.

Only pole terms contribute to $M_{9,K}$, since polynomial moments do
not depend on $X$.  Cancel one factor $D_N$ from $D_N^6/D_K$ and take
the residues at $-j^2$, $N<j\leq K$.  With $u_j=j/K$, source (C.3)
becomes
\begin{equation}
 \begin{aligned}
 \mathbf q^TM_{9,K}\mathbf q
    &=A_K\sum_{j=N+1}^{K}r_jp(-u_j^2)^2,\\
 r_j&=\frac{u_j^8\displaystyle\prod_{\ell=1}^{N}
                      (u_\ell^2-u_j^2)^5}
                 {\displaystyle\prod_{\substack{N<k\leq K\\k\ne j}}
                      (u_k^2-u_j^2)}.
 \end{aligned}
 \label{c-residues}
\end{equation}
For clarity, before scaling the residue multiplier is
$j^8D_N(-j^2)^5/\prod_{N<k\leq K,\,k\ne j}(k^2-j^2)$.
Its power of $K$ is
$8+10N-2(h-1)=12N-2K+10$, confirming the same $A_K$ as in the integral.
The factor $u_j^8$ is the sole change in this scaled residue formula
from the fourth power appearing in the source.
The numerator has absolute value at most one.  For each denominator
factor, $u_k+u_j\geq2\alpha$ and
$|u_k-u_j|=|k-j|/K$.  Setting $m=h-1$ gives the analogue of source (C.4):
\begin{equation}
 \begin{aligned}
 |r_j|&\leq
 \frac{(K/(2\alpha))^{h-1}}{(j-N-1)!(K-j)!}
 \leq\frac{(K/\alpha)^{h-1}}{(h-1)!}\\
 &\leq\exp\{\lambda K(1-\log(\alpha\lambda))\}.
 \end{aligned}
 \label{c-residue-bound}
\end{equation}
The middle step uses $\binom m{j-N-1}\leq2^m$.  For the last step,
$m!\geq(m/e)^m$ and the function $x\log(eK/(\alpha x))$ is increasing
on $0<x\leq h$.

Expand $p$ in the Legendre polynomials for $[a,b]$.  At $-1\leq t\leq0$
their transformed argument has absolute value at most
$Y=(2+a+b)/(b-a)$.  The source's Legendre bound $|P_j(s)|\leq R^j$,
where $R=Y+\sqrt{Y^2-1}$, and Cauchy--Schwarz give
\begin{equation}
 \max_{-1\leq t\leq0}|p(t)|^2\leq
 \frac{h^2R^{2h}}{b-a}\int_a^b p(t)^2\,dt,
 \qquad Y=\frac{2+a+b}{b-a},\quad R=Y+\sqrt{Y^2-1}.
 \label{c-legendre}
\end{equation}
Indeed, the squared norms are $(b-a)/(2j+1)$ and
$\sum_{j=0}^{h-1}(2j+1)=h^2$.  This is exactly the step in source (C.5).
There are $h$ residues, so \eqref{c-gram-lower}--\eqref{c-legendre} imply
\begin{equation}
 \begin{aligned}
 \|B_K\|&=
  \sup_{\mathbf q\ne0}
  \frac{|\mathbf q^TM_{9,K}\mathbf q|}{\mathbf q^TG\mathbf q}
  \leq\frac{h^3}{c_a(b-a)}e^{\kappa(a,b)K},\\
 \kappa(a,b)&=F(a,b)+\lambda(1-\log(\alpha\lambda))+2\lambda\log R.
 \end{aligned}
 \label{c-kappa-bound}
\end{equation}
For the source choice $a=1/18$, $b=1/3$, we have $Y=43/5$ and
$R=(43+4\sqrt{114})/5$.  Thus source (C.6) is unchanged:
\begin{equation}
 \begin{aligned}
 \kappa={}&\frac{2\pi}{\sqrt3}+\log\frac43
       +\frac9{20}\log18
       +\frac{37}{40}\left(1+\log\frac{1600}{111}\right)
       +\frac{37}{20}\log\frac{43+4\sqrt{114}}5\\
   &\simeq13.865856678865<\frac{6933}{500}.
 \end{aligned}
 \label{c-kappa-value}
\end{equation}
\noindent\textbf{[Conditional]} Using the elementary enclosure from
source Appendix~A for this strict numerical inequality, its positive
margin absorbs $h^3/(c_a(b-a))$ for all sufficiently large $K$, proving
\eqref{c-norm-bound}.  The exponent is unchanged; the threshold need
not be the source's threshold.
More precisely, the retained enclosure is
$\kappa<6932928339433/500000000000<6933/500$ from source (A.11).
Absorption uses only $\log(h^3/(c_a(b-a)))=O(\log K)$, whereas the
positive gap between the two exponents is multiplied by $K$.

\subsection{Perturbation and the conditional approximation statement}
\label{c-approximation-conditional}

Fix a cutoff $M$ and abbreviate $Q_K=Q_{9,K,M}$.  Affine dependence
$G_{9,K}(\xi+\delta)=G+\delta M_{9,K}$ gives the exact identity
\begin{equation}
 Q_K(\xi+\delta)=Q_K(\xi)\det(I+\delta B_K)
      =Q_K(\xi)\prod_{j=1}^{h}(1+\delta\beta_j),
 \label{c-determinant-perturbation}
\end{equation}
where $\beta_j$ are the real eigenvalues of $B_K$.  If
$|\delta|\|B_K\|<1$, every factor is positive.  Moreover
\[
 \det(I+\delta B_K)
 \leq(1+|\delta|\|B_K\|)^h
 \leq\exp\{h|\delta|\|B_K\|\}.
\]
This explains how a polynomial that is small at $\xi$ can be kept
nonzero at a sufficiently close rational number.
The scalar $m_{9,K,M}S_K$ cancels in this relative comparison; it is
independent of $X$.  Thus the analytic radius of control comes from
$B_K$, while the arithmetic normalization is needed later to turn the
rational evaluation into an integer.  These are separate uses of the
same determinant identity.

\noindent\textbf{[Conditional]} To continue as in source (C.7), assume
the additional input
\begin{equation}
 0<Q_K(\xi)\leq e^{-\varepsilon K^2},\qquad \varepsilon>0,
 \quad\text{for all sufficiently large allowed }K.
 \label{c-negative-hypothesis}
\end{equation}
Choose constants $c>\lambda/\varepsilon$ and $\nu>\rho c$.  For a
sufficiently large positive integer $b$, take
$K=40\lceil c\log b/40\rceil$, so $c\log b\leq K<c\log b+40$.
If $|a/b-\xi|\leq b^{-\nu}$, put $\delta=a/b-\xi$.  The relative norm
bound gives
\[
 |\delta|\|B_K\|\leq e^{40\rho}b^{\rho c-\nu}=o(K^{-1}).
\]
The determinant in \eqref{c-determinant-perturbation} is consequently
positive and at most two eventually, uniformly in $a$.  Integrality
and degree $h$ would then force
\begin{equation}
 \begin{aligned}
 0<b^hQ_K(a/b)
 &\leq2\exp\{\lambda K\log b-\varepsilon K^2\}\\
 &\leq2\exp\{-(\varepsilon-\lambda/c)K^2\}<1,
 \qquad b^hQ_K(a/b)\in\mathbb Z,
 \end{aligned}
 \label{c-rational-contradiction}
\end{equation}
an impossibility.  Under the stated hypotheses this excludes such
approximations for every $\nu>\rho\lambda/\varepsilon$, after choosing
$c$ between $\lambda/\varepsilon$ and $\nu/\rho$.

\noindent\textbf{[Blocked conclusion]} Section~7 supplies
$B_{9,100000}\simeq1.974551870160>0$, not the negative hypothesis
\eqref{c-negative-hypothesis}.  The source's choices
$\varepsilon=79/1600$, $c=75/4$, and $\nu=260$ therefore do not follow
for this ninth-order construction.  The calculation above describes
the consequence of a future negative budget, not an approximation
bound obtained here.

\subsection{Coefficient height and primitive normalization}
\label{c-height-content}

\noindent\textbf{[Derived]} For $Q\in\mathbb R[X]$, let $H(Q)$ be the
largest absolute value of its coefficients.  Since
$1<\xi<2$ and \eqref{c-determinant-perturbation} is a polynomial
identity, the sum of the absolute coefficient values is bounded by
\[
 H(Q_K)\leq Q_K(\xi)
       \prod_{j=1}^{h}\bigl(1+(1+|\xi|)|\beta_j|\bigr)
 \leq Q_K(\xi)(1+3e^{\rho K})^h.
\]
\noindent\textbf{[Conditional]} With the arithmetic and real estimates
already stated, this gives the analogue of source (C.8), without
requiring any additional negative-budget assumption:
\begin{equation}
 \begin{aligned}
 \log H(Q_{9,K,M})
     &\leq(B_{9,M}+\lambda\rho)K^2+o_M(K^2),
       &\lambda\rho&=\frac{256521}{20000}=12.82605,\\
 \log H(Q_K)&\leq(\lambda\rho-\varepsilon)K^2+o(K^2)
       &&\text{if \eqref{c-negative-hypothesis} holds}.
 \end{aligned}
 \label{c-height-bound}
\end{equation}
For $M=100000$ the coefficient in the first line is approximately
$14.800601870160$.  Thus a height estimate remains available, although
the source's height coefficient based on its negative margin is not
transferred.

For nonzero $A\in\mathbb Q[X]$, define $\operatorname{cont}_{\mathbb Q}(A)>0$
so that $A/\operatorname{cont}_{\mathbb Q}(A)$ is an integer polynomial
whose coefficients have greatest common divisor one.  Put
\begin{equation}
 P_K=\frac{\Delta_{9,K}}{\operatorname{cont}_{\mathbb Q}(\Delta_{9,K})},
 \qquad Q_{9,K,M}=a_{K,M}P_K,\qquad a_{K,M}\in\mathbb Z_{>0}.
 \label{c-primitive-normalization}
\end{equation}
The last assertion requires the integrality of $Q_{9,K,M}$ from
Section~5.  To see it, the multiplier is initially a positive rational
number.  A denominator in lowest terms would have to divide every
coefficient of the primitive polynomial $P_K$, so that denominator is
one.  Equivalently, a Bezout combination of its coefficients equals
one and forces the multiplier to be integral.  Hence
\begin{equation}
 0<P_K(\xi)=\frac{Q_{9,K,M}(\xi)}{a_{K,M}}
       \leq Q_{9,K,M}(\xi),\qquad
 H(P_K)=\frac{H(Q_{9,K,M})}{a_{K,M}}\leq H(Q_{9,K,M}).
 \label{c-primitive-bounds}
\end{equation}
The primitive polynomial is independent of the cutoff $M$.  All
available upper bounds descend to it, including any hypothetical
small-value bound.  This removal of content does not itself produce a
negative quadratic coefficient; no size estimate for $a_{K,M}$ has
been obtained that would supply the missing decay.
